\section{The three main theorems}

The three statements below are proved with their existing hypotheses and model choices. Theorem B means the full statement with literal simplicial structure, not only its earlier seven-clause core.

\begin{resultblock}
\begin{theorem}[Theorem A: the homotopically constant resolution model]
\label{catalogue:MainA}
There is a left proper combinatorial model structure
$\sSpace^{\dag\vee}_{\CSS}:=L_{S_h}\sSpace^{\dag\vee}_{\Reedy}$.
Its cofibrations are exactly the Reedy cofibrations and its weak equivalences are exactly the $S_h$-local equivalences.
Its fibrant objects are precisely the Reedy fibrant diagrams $W$ for which the canonical counit $\ell_W:cW_0\to W$ is a levelwise dagger Joyal weak equivalence.
The constant--evaluation adjunction is a Quillen equivalence
\begin{align}
 c:\sSet^{\dag}_{\Joyal}\rightleftarrows\sSpace^{\dag\vee}_{\CSS}:\ev_0.
\end{align}
The same existential model structure satisfies all these clauses.
\end{theorem}

\Source{DaggerModels.PartII.theoremA;
DaggerModelsPartIICompanion.mainA}{lean/DaggerModelsPartII/ProofTheoremA.lean;
lean/MainTheoremsPartII.lean}
\CatalogueNote{Scope}{Arbitrary original value-universe parameter u.}
\CatalogueNote{Manuscript context}{Theorem A, cs.def.Sh, cs.thm.resolution, cs.thm.equivalence}
\CatalogueNote{Construction and reuse}{Existing original theorem and public alias; original hypotheses and statement retained.}
\end{resultblock}

\begin{resultblock}
\begin{theorem}[Theorem B: dagger complete Segal models and change of index]
\label{catalogue:MainB}
The injective localization $L_{S_{\rev}}\sSpace^{\dag}_{\rev,\inj}$ is left proper, combinatorial and simplicial for the literal mapping spaces and tensor.
Its fibrant objects are exactly the injectively fibrant unitarily complete dagger Segal spaces.
The projective localization has the levelwise Kan diagrams satisfying the corresponding derived Segal and unitary completeness conditions as its fibrant objects.
The actual-index projective localization has precisely the levelwise Kan diagrams whose restriction along $u$ satisfies that derived condition.
The actual adjunction
\begin{align}
 u_!:L_{\widetilde S_{\rev}}\sSpace^{\dag}_{\rev,\proj}
 \rightleftarrows L_{u_!(\widetilde S_{\rev})}\sSpace^{\dag}_{\proj}:u^*
\end{align}
is a Quillen equivalence.
The full Lean theorem also supplies the native simplicial enrichment, its natural comparison with the literal Hom bifunctor, ordinary-morphism normalization, parameterized tensor and cotensor adjunctions, and natural exponential laws.
\end{theorem}

\Source{DaggerModels.PartII.theoremB;
DaggerModelsPartIICompanion.mainB}{lean/DaggerModelsPartII/ProofTheoremB.lean;
lean/MainTheoremsPartII.lean}
\CatalogueNote{Scope}{The assembled theorem, including its actual localized models and literal simplicial structure, has value universe zero.}
\CatalogueNote{Manuscript context}{Theorem B, h1.lem.rev-models, h1.prop.fibrant, h1.thm.main}
\CatalogueNote{Construction and reuse}{Existing original theorem and public alias; original hypotheses and statement retained.}
\end{resultblock}

\begin{resultblock}
\begin{theorem}[Theorem C: the original four-link comparison]
\label{catalogue:MainC}
The following four adjunctions are Quillen equivalences.
Each line displays its left adjoint first, with the exact original models and directions:
\begin{align}
 \frakC_{\dag}:\sSet^{\dag}_{\Joyal}
 &\rightleftarrows\sCat^{\dag}_{\Bergner}:\N_{\dag},\\
 \delta:\sSet^{\dag}_{\Joyal}
 &\rightleftarrows L_{S_{\rev}}\sSpace^{\dag}_{\rev,\inj}:i_0^*,\\
 \mathrm{id}:L_{\widetilde S_{\rev}}\sSpace^{\dag}_{\rev,\proj}
 &\rightleftarrows L_{S_{\rev}}\sSpace^{\dag}_{\rev,\inj}:\mathrm{id},\\
 u_!:L_{\widetilde S_{\rev}}\sSpace^{\dag}_{\rev,\proj}
 &\rightleftarrows L_{u_!(\widetilde S_{\rev})}\sSpace^{\dag}_{\proj}:u^*.
\end{align}
In the manuscript's left-to-right zigzag, the first and third rightward arrows instead display the right adjoints.
The original fixed Lean target has no additional assumptions;
in particular ordinary completeness of the original input is not imposed.
\end{theorem}

\Source{DaggerModels.PartII.TheoremCOriginalCompletionProof.original_statement;
DaggerModelsPartIICompanion.mainC;
DaggerModels.PartII.TheoremCOriginalFourLinksAudit.partI;
DaggerModels.PartII.TheoremCOriginalFourLinksAudit.constantVertex;
DaggerModels.PartII.TheoremCOriginalFourLinksAudit.reversalIdentity;
DaggerModels.PartII.TheoremCOriginalFourLinksAudit.indexChange}{lean/TheoremCStages/TheoremCOriginalCompletionProof.lean;
lean/MainTheoremCPartII.lean;
lean/TheoremCStages/TheoremCOriginalFourLinksAudit.lean}
\CatalogueNote{Scope}{Original four-link statement at value universe zero.}
\CatalogueNote{Manuscript context}{Theorem C, jt.thm.main, jt.cor.goal}
\CatalogueNote{Construction and reuse}{Existing original theorem and public alias; original hypotheses and statement retained.}
\end{resultblock}
